% =====================================================================
% Capítulo 2 — Marco teórico
% =====================================================================
\chapter{Marco teórico: fundamentos de los modelos de difusión}
\label{cap:marco-teorico}

Este capítulo concentra todo el aparato teórico del trabajo. Primero
se presenta la idea de los modelos de difusión y sus dos raíces
históricas; después se desarrolla la formulación matemática completa
—proceso hacia delante, objetivo de entrenamiento, parametrizaciones
de la predicción, calendarios de ruido, muestreadores y guiado para
problemas inversos—; a continuación se repasan brevemente sus
aplicaciones y los fundamentos físicos del escaneo PET; y se cierra
con la revisión del estado del arte en reconstrucción de PET con
difusión. Los capítulos posteriores \emph{referencian} las ecuaciones
aquí definidas sin volver a derivarlas.

% ---------------------------------------------------------------------
\section{¿Qué es un modelo de difusión?}
\label{sec:que-es-difusion}

Un \emph{modelo de difusión} es un modelo generativo probabilístico
que aprende a producir muestras nuevas de una distribución de datos
$p(x)$ invirtiendo, paso a paso, un proceso de corrupción gradual.
La idea central es: si se sabe cómo destruir la estructura de un dato 
añadiéndole ruido de forma progresiva, hasta
convertirlo en ruido puro, entonces aprender a \emph{deshacer} este proceso 
equivale a aprender a \emph{crear} datos a partir de
ruido. El modelo se entrena, no para describir
explícitamente la distribución $p(x)$ —algo intratable para datos de
alta dimensión como las imágenes— sino para revertir un proceso de
difusión conocido.

Conviene fijar desde el principio dos conceptos elementales que
reaparecerán a lo largo de toda la sección:

\begin{itemize}
    \item El \textbf{proceso hacia delante} (\textit{forward} o
    \textit{diffusion process}): una cadena fija, sin parámetros
    entrenables, que añade ruido gaussiano a un dato $x_0$ a lo largo
    de $T$ pasos hasta que $x_T$ es indistinguible de una muestra de
    $\mathcal{N}(0,\mathbf{I})$.
    \item El \textbf{proceso hacia atrás} (\textit{reverse} o
    \textit{denoising process}): la cadena que queremos aprender, en
    la que una red neuronal con parámetros $\theta$ estima cómo
    eliminar el ruido en cada paso para reconstruir progresivamente
    una muestra de datos a partir de ruido puro.
\end{itemize}

A diferencia de los GAN, donde un generador y un discriminador se
entrenan de forma adversaria, o de los VAE, que comprimen los datos
en un espacio latente de baja dimensión, los modelos de difusión
descomponen la generación en una secuencia larga de pasos de
\emph{denoising} pequeños y bien condicionados, cada uno de los cuales
es un problema de aprendizaje supervisado sencillo y estable. 

% ---------------------------------------------------------------------
\section{Orígenes teóricos}
\label{sec:por-que-difusion}

\subsection{El origen físico: termodinámica del no equilibrio}
\label{subsec:origen-fisico}

La motivación original de los modelos de difusión es explícitamente
física. El trabajo fundacional de Sohl-Dickstein et
al.~\cite{sohldickstein2015} se inspira en la \emph{termodinámica del
no equilibrio}: del mismo modo que una gota de tinta en un vaso de
agua se difunde hasta repartirse de manera homogénea —un estado de
máxima entropía—, una distribución de datos
compleja puede transformarse, mediante un proceso de difusión, en una
distribución gaussiana simple. La apuesta teórica fue que, si dicho
proceso se aplica con pasos suficientemente pequeños, cada paso de la
\emph{inversión} admite también una forma funcional sencilla
(gaussiana), lo que lo hace aprendible por una red neuronal. Es
decir: se renuncia a modelar directamente la complejísima $p(x)$ y, a
cambio, se modela una larga secuencia de transiciones locales, cada
una de ellas tratable.

\subsection{Score-matching: la otra raíz teórica}
\label{subsec:score-matching}

Existe una segunda línea de trabajo, inicialmente independiente: el
\emph{score-matching}. En esta formulación, propuesta por Song y
Ermon~\cite{song2019scorematching}, el objetivo no es modelar la
densidad $p(x)$ sino su \emph{score}, el gradiente del logaritmo de
la densidad respecto a los datos, $\nabla_x \log p(x)$. Conocer este
campo vectorial basta para generar muestras: partiendo de ruido, se
siguen iterativamente los gradientes hacia regiones de alta densidad
mediante \emph{dinámica de Langevin}. El problema práctico es que el
score está mal estimado en las zonas de baja densidad, precisamente
donde el muestreo arranca; la solución que propusieron —perturbar los
datos con ruido gaussiano a múltiples escalas y estimar el score de
cada versión ruidosa— es, en la práctica, casi idéntica al proceso de
difusión de Sohl-Dickstein. Esta convergencia se formalizó en el
marco unificado de las ecuaciones diferenciales estocásticas de Song
et al.~\cite{song2021sde}, donde tanto los DDPM como los modelos de
score-matching se revelan como dos discretizaciones de la misma SDE
subyacente.

\subsection{Perspectiva de nuestro estudio}
\label{subsec:aplica-hoy}

La intuición física de la difusión y el marco de score-matching son
dos perspectivas teóricas distintas sobre el mismo objeto, y ambas
describen correctamente lo que hacen los modelos modernos. Sin
embargo, la arquitectura con la que se entrena e infiere en
este trabajo es la de la familia \textbf{DDPM}~\cite{ho2020ddpm} y su
muestreador determinista \textbf{DDIM}~\cite{song2021ddim}: el
entrenamiento no se plantea como score-matching ni como la inversión
de una SDE, sino como un objetivo de \emph{predicción del ruido} (o de
una reparametrización suya, como se verá en la
\Cref{subsec:parametrizaciones}). No obstante, ese objetivo realiza
\emph{implícitamente} una estimación del score: la
\Cref{eq:score-eps} muestra que ambas cantidades coinciden salvo un
factor de reescalado dependiente de $t$. Esta equivalencia es la que
permitirá, en la \Cref{subsec:difusion-condicional}, guiar un modelo
no condicional con una observación externa usando razonamiento
bayesiano sobre scores.

% ---------------------------------------------------------------------
\section{Formulación matemática}
\label{sec:formulacion-difusion}

En el contexto de generación de datos, se parte de una distribución
$p(x)$ que representa los datos de entrenamiento —por ejemplo, todas
las imágenes PET de dosis completa—. Esta distribución es demasiado
compleja para admitir una expresión cerrada, pero generar imágenes
nuevas equivale a muestrear de ella. Los modelos de difusión
resuelven el problema construyendo un proceso que transforma
cualquier distribución en una gaussiana estándar, y aprendiendo a
invertirlo.

\subsection{Proceso hacia delante}
\label{subsec:forward}

El proceso hacia delante es una cadena de Markov: cada paso de
adición de ruido depende exclusivamente del paso anterior. La
probabilidad conjunta de la trayectoria completa, y su contraparte
aprendida hacia atrás, se factorizan como
\begin{align}
    q(x_{1:T} \mid x_0)
        &= \prod_{t=1}^{T} q(x_t \mid x_{t-1}),
        \label{eq:forward-chain}\\
    p_\theta(x_{0:T})
        &= p_\theta(x_T) \prod_{t=1}^{T} p_\theta(x_{t-1} \mid x_t).
        \label{eq:reverse-chain}
\end{align}

Cada transición hacia delante añade una cantidad controlada de ruido
gaussiano:
\begin{equation}
    q(\mathbf{x}_t \mid \mathbf{x}_{t-1})
    = \mathcal{N}\!\left(\mathbf{x}_t;\, \sqrt{1 - \beta_t}\,
      \mathbf{x}_{t-1},\, \beta_t \mathbf{I}\right),
    \label{eq:forward-step}
\end{equation}
donde $\{\beta_t\}_{t=1}^{T}$ es un calendario de varianza
predefinido (\Cref{subsec:schedules}). Definiendo
$\alpha_t = 1 - \beta_t$ y $\bar{\alpha}_t = \prod_{s=1}^{t}
\alpha_s$, la composición de transiciones gaussianas admite una forma
cerrada que permite muestrear $\mathbf{x}_t$ directamente desde
$\mathbf{x}_0$ en un único paso:
\begin{equation}
    q(\mathbf{x}_t \mid \mathbf{x}_0)
    = \mathcal{N}\!\left(\mathbf{x}_t;\, \sqrt{\bar{\alpha}_t}\,
      \mathbf{x}_0,\, (1 - \bar{\alpha}_t)\mathbf{I}\right),
    \label{eq:forward-closed}
\end{equation}
o, por reparametrización,
\begin{equation}
    \mathbf{x}_t
    = \sqrt{\bar{\alpha}_t}\,\mathbf{x}_0
    + \sqrt{1 - \bar{\alpha}_t}\,\boldsymbol{\varepsilon},
    \qquad
    \boldsymbol{\varepsilon} \sim \mathcal{N}(\mathbf{0}, \mathbf{I}).
    \label{eq:forward-reparam}
\end{equation}
Esta forma cerrada es la que hace eficiente el entrenamiento: en cada
iteración se puede construir una muestra ruidosa de cualquier nivel
$t$ sin simular la cadena completa.

\subsection{Proceso inverso y objetivo de entrenamiento}
\label{subsec:objetivo}

El proceso inverso se modela con transiciones gaussianas cuya media
estima una red neuronal:
\begin{equation}
    p_\theta(\mathbf{x}_{t-1} \mid \mathbf{x}_t)
    = \mathcal{N}\!\left(\mathbf{x}_{t-1};\,
      \boldsymbol{\mu}_\theta(\mathbf{x}_t, t),\,
      \sigma_t^2 \mathbf{I}\right).
    \label{eq:reverse-step}
\end{equation}
El entrenamiento maximiza la verosimilitud de los datos a través de
una cota variacional (ELBO). Ho et al.~\cite{ho2020ddpm} demostraron que, 
tras simplificar los pesos de los términos de la cota, el objetivo se 
reduce a un problema de regresión notablemente simple: predecir el ruido
$\boldsymbol{\varepsilon}$ que se añadió en la
\Cref{eq:forward-reparam},
\begin{equation}
    \mathcal{L}_{\text{simple}}
    = \mathbb{E}_{t \sim \mathcal{U}[1,T],\, \mathbf{x}_0,\,
                   \boldsymbol{\varepsilon}}
      \!\left[\left\|\boldsymbol{\varepsilon}
        - \boldsymbol{\varepsilon}_\theta\!\left(\mathbf{x}_t, t\right)
        \right\|^2\right].
    \label{eq:loss-eps}
\end{equation}
En cada iteración se muestrea un dato $\mathbf{x}_0$, un paso $t$
uniforme y un ruido $\boldsymbol{\varepsilon}$; se construye
$\mathbf{x}_t$ con la \Cref{eq:forward-reparam}; se predice el ruido como
veremos posteriormente en la \Cref{subsec:parametrizaciones}; y se actualizan los
parámetros con el gradiente del error cuadrático medio sobre el ruido. Cada paso es
un problema supervisado estable, sin dinámica adversaria.

\subsection{Parametrizaciones de la predicción}
\label{subsec:parametrizaciones}

La \Cref{eq:forward-reparam} liga linealmente tres cantidades
—$\mathbf{x}_t$, $\mathbf{x}_0$ y $\boldsymbol{\varepsilon}$—, de
modo que conocidas dos de ellas la tercera queda determinada. Esto da
libertad para elegir \emph{qué} predice la red, y la elección tiene
consecuencias prácticas:

\paragraph{Predicción de ruido ($\epsilon$-prediction).}
Es la formulación original de la \Cref{eq:loss-eps}. Conocido el
ruido estimado $\hat{\boldsymbol{\varepsilon}}$, la imagen limpia se
recupera invirtiendo la \Cref{eq:forward-reparam}:
\begin{equation}
    \hat{\mathbf{x}}_0
    = \frac{\mathbf{x}_t
        - \sqrt{1-\bar{\alpha}_t}\;\hat{\boldsymbol{\varepsilon}}}
      {\sqrt{\bar{\alpha}_t}}.
    \label{eq:x0-from-eps}
\end{equation}
Su punto débil está en los \textit{timesteps} extremos: cuando
$t \to T$, $\sqrt{\bar{\alpha}_t} \to 0$ y la
\Cref{eq:x0-from-eps} amplifica enormemente cualquier error de la
red.

\paragraph{Predicción de la imagen ($\mathbf{x}_0$-prediction).}
La red estima directamente $\mathbf{x}_0$. Tiene el problema
simétrico: en los \textit{timesteps} de poco ruido ($t \to 0$) la
tarea es casi trivial y la señal de aprendizaje sobre el ruido
residual se desvanece.

\paragraph{Predicción de velocidad (\textit{v-prediction}).}
Salimans y Ho~\cite{salimans2022vprediction} proponen el objetivo
intermedio
\begin{equation}
    \mathbf{v}_t
    = \sqrt{\bar{\alpha}_t}\,\boldsymbol{\varepsilon}
    - \sqrt{1-\bar{\alpha}_t}\,\mathbf{x}_0,
    \label{eq:v-def}
\end{equation}
una combinación que rota suavemente entre predecir el ruido (en
$t \to 0$) y predecir la imagen con signo cambiado (en $t \to T$).
La pérdida correspondiente es
\begin{equation}
    \mathcal{L}_{\mathbf{v}}
    = \mathbb{E}_{t,\, \mathbf{x}_0,\, \boldsymbol{\varepsilon}}
      \!\left[\left\|\mathbf{v}_t
        - \mathbf{v}_\theta(\mathbf{x}_t,\, t)
      \right\|^2\right],
    \label{eq:loss-v}
\end{equation}
y la estimación de la imagen limpia a partir de la predicción
$\hat{\mathbf{v}}_t$ es
\begin{equation}
    \hat{\mathbf{x}}_0
    = \sqrt{\bar{\alpha}_t}\,\mathbf{x}_t
    - \sqrt{1-\bar{\alpha}_t}\,\hat{\mathbf{v}}_t.
    \label{eq:x0-from-v}
\end{equation}
La parametrización $\mathbf{v}$ es numéricamente estable en los dos
límites $t \to 0$ y $t \to T$, lo que evita las inestabilidades de
entrenamiento de las dos anteriores en los \textit{timesteps} de muy
bajo o muy alto ruido. Es la parametrización empleada en los
\textit{pipelines} de reconstrucción PET de este trabajo
(\Cref{sec:pet-entrenamiento}); el experimento de mariposas
(\Cref{cap:mariposas}) emplea la $\epsilon$-prediction original, lo
que permite comparar ambas en la práctica.

\subsection{Conexión con el score}
\label{subsec:score-conexion}

Aunque los DDPM no se derivan del score-matching, sí lo realizan
implícitamente. Para la distribución marginal ruidosa
$q(\mathbf{x}_t)$ inducida por la \Cref{eq:forward-closed} se cumple
\begin{equation}
    \nabla_{\mathbf{x}_t} \log q(\mathbf{x}_t)
    = -\frac{\boldsymbol{\varepsilon}_\theta(\mathbf{x}_t, t)}
            {\sqrt{1-\bar{\alpha}_t}},
    \label{eq:score-eps}
\end{equation}
es decir, la red que predice el ruido es —salvo un factor de
reescalado dependiente de $t$— un estimador del score de la
distribución ruidosa. Las parametrizaciones de la
\Cref{subsec:parametrizaciones} son interconvertibles, así que la
identidad vale para cualquiera de ellas. Esta equivalencia es la
puerta de entrada del razonamiento bayesiano sobre el muestreo que
fundamenta el guiado posterior
(\Cref{subsec:difusion-condicional}).

\subsection{Calendarios de ruido}
\label{subsec:schedules}

El calendario $\{\beta_t\}$ determina cuánta señal sobrevive en cada
\textit{timestep} —la relación señal-ruido es
$\mathrm{SNR}(t) = \bar{\alpha}_t / (1-\bar{\alpha}_t)$— y, con ello,
sobre qué niveles de ruido se concentra el esfuerzo de aprendizaje.
El calendario \textbf{lineal} original de Ho et al.~\cite{ho2020ddpm}
interpola $\beta_t$ entre $10^{-4}$ y $0{,}02$; destruye la señal
demasiado deprisa, de modo que muchos \textit{timesteps} finales
resultan casi inútiles para el aprendizaje. Nichol y
Dhariwal~\cite{nichol2021cosine} proponen el calendario
\textbf{cosenoidal}, definido sobre la señal acumulada:
\begin{equation}
    \bar{\alpha}_t = \frac{f(t)}{f(0)},
    \qquad
    f(t) = \cos^2\!\left(
        \frac{t/T + s}{1 + s}\cdot\frac{\pi}{2}
    \right),
    \label{eq:cosine-schedule}
\end{equation}
con un pequeño \textit{offset} $s$ que evita degeneración en $t=0$.
Este perfil degrada la señal de forma más gradual y produce una
distribución de SNR mejor adaptada a resoluciones medias; es el
calendario empleado en los \textit{pipelines} PET de este trabajo.

\subsection{Muestreo: DDPM y DDIM}
\label{subsec:muestreo}

\paragraph{Muestreo ancestral (DDPM).}
Para generar, se parte de ruido puro
$\mathbf{x}_T \sim \mathcal{N}(\mathbf{0}, \mathbf{I})$ y se aplica
$T$ veces la transición inversa aprendida. En la parametrización
$\epsilon$, cada paso recupera
\begin{equation}
    \mathbf{x}_{t-1}
    = \frac{1}{\sqrt{\alpha_t}}
      \!\left(\mathbf{x}_t
        - \frac{\beta_t}{\sqrt{1 - \bar{\alpha}_t}}\,
        \boldsymbol{\varepsilon}_\theta(\mathbf{x}_t, t)\right)
    + \sigma_t\,\mathbf{z},
    \qquad
    \mathbf{z} \sim \mathcal{N}(\mathbf{0}, \mathbf{I}).
    \label{eq:ddpm-step}
\end{equation}
El procedimiento es estocástico y requiere los $T$ pasos completos
(típicamente $T=1000$), lo que lo hace lento en inferencia.

\paragraph{Muestreo determinista acelerado (DDIM).}
Song et al.~\cite{song2021ddim} observan que existe una familia de
procesos no markovianos con las mismas marginales que la
\Cref{eq:forward-closed}, lo que permite saltar entre
\textit{timesteps} no consecutivos. Cada paso DDIM se escribe en
función de la estimación de la imagen limpia $\hat{\mathbf{x}}_0$
(obtenida con la \Cref{eq:x0-from-eps} o la \Cref{eq:x0-from-v},
según la parametrización):
\begin{equation}
    \mathbf{x}_{t-1}
    = \sqrt{\bar{\alpha}_{t-1}}\,\hat{\mathbf{x}}_0
    + \sqrt{1-\bar{\alpha}_{t-1}-\tilde{\sigma}_t^2}\,
      \hat{\boldsymbol{\varepsilon}}
    + \tilde{\sigma}_t\,\mathbf{z},
    \label{eq:ddim-step}
\end{equation}
donde el parámetro $\eta$ escala $\tilde{\sigma}_t$ entre el caso
DDPM ($\eta=1$) y el caso totalmente \textbf{determinista}
($\eta=0$), en el que la trayectoria queda fijada por la semilla
inicial. En la práctica, DDIM con $\eta=0$ y 50 pasos produce
muestras de calidad comparable a los 1000 pasos ancestrales con una
veinteava parte del coste; es el muestreador empleado en la
inferencia de los \textit{pipelines} PET.

\subsection{Difusión condicional y problemas inversos}
\label{subsec:difusion-condicional}

Hasta aquí el modelo genera muestras incondicionales de $p(x)$. En
un problema de reconstrucción interesa la distribución
\emph{posterior} $p(\mathbf{x} \mid \mathbf{y})$: qué imágenes
limpias son compatibles con una observación degradada $\mathbf{y}$.
Existen dos estrategias.

\paragraph{Condicionamiento supervisado.}
Si se dispone de pares $(\mathbf{x}, \mathbf{y})$, la vía directa es
entrenar la red de \textit{denoising} con la observación como entrada
adicional, $\mathbf{v}_\theta(\mathbf{x}_t, \mathbf{y}, t)$
—típicamente concatenando $\mathbf{y}$ como canal extra—. El
objetivo de entrenamiento es el mismo de la \Cref{eq:loss-v}, y el
muestreo (\Cref{eq:ddim-step}) hereda el condicionamiento sin
cambios. El modelo aprende directamente
$p(\mathbf{x} \mid \mathbf{y})$ y aprovecha íntegramente los datos
pareados, a costa de quedar ligado a la degradación concreta vista en
entrenamiento.

\paragraph{Guiado posterior (DPS).}
La alternativa es entrenar un prior \emph{no condicional}
$p(\mathbf{x})$ e inyectar la observación solo en inferencia. El
punto de partida es la regla de Bayes aplicada al score de la
distribución ruidosa:
\begin{equation}
    \nabla_{\mathbf{x}_t} \log p(\mathbf{x}_t \mid \mathbf{y})
    = \nabla_{\mathbf{x}_t} \log p(\mathbf{x}_t)
    + \nabla_{\mathbf{x}_t} \log p(\mathbf{y} \mid \mathbf{x}_t).
    \label{eq:posterior-score}
\end{equation}
El primer término lo proporciona el modelo entrenado
(\Cref{eq:score-eps}); el segundo es intratable porque relaciona la
observación con el estado \emph{ruidoso} $\mathbf{x}_t$.
\textit{Diffusion Posterior Sampling} (DPS)~\cite{chung2023dps}
propone la aproximación
$p(\mathbf{y} \mid \mathbf{x}_t) \approx
p(\mathbf{y} \mid \hat{\mathbf{x}}_0(\mathbf{x}_t))$, usando la
estimación de imagen limpia que el propio modelo produce en cada paso
(\Cref{eq:x0-from-v}). Para un modelo de observación gaussiano
$\mathbf{y} = A\mathbf{x} + \mathbf{n}$, el término de verosimilitud
se convierte en un paso de gradiente sobre el error de medición, que
se intercala en cada iteración del muestreo:
\begin{equation}
    \mathbf{x}_t \leftarrow
        \mathbf{x}_t - \omega \cdot
        \nabla_{\mathbf{x}_t}
        \!\left[\frac{1}{2}
            \left\|A\,\hat{\mathbf{x}}_0(\mathbf{x}_t)
              - \mathbf{y}\right\|_2^2
        \right],
    \label{eq:dps-update}
\end{equation}
donde $\omega$ es la escala de guiado, único hiperparámetro del
procedimiento: valores bajos confían en el prior y valores altos
fuerzan la consistencia con la medición. El gradiente atraviesa la
red ($\hat{\mathbf{x}}_0$ depende de $\theta$), por lo que cada paso
de inferencia requiere una pasada hacia atrás además de la pasada
hacia delante. La gran ventaja del enfoque es que el mismo prior
sirve para cualquier operador $A$ sin reentrenar; su debilidad, que
la calidad depende de lo bien que el modelo de observación describa
la degradación real.

% ---------------------------------------------------------------------
\section{Principales aplicaciones}
\label{sec:aplicaciones-difusion}

El interés teórico se vio rápidamente respaldado por resultados
empíricos que situaron a la difusión por delante de los GAN en
síntesis de imágenes~\cite{dhariwal2021beatgans}. Desde entonces, el
abanico de aplicaciones se ha ampliado de forma notable. A modo de
panorámica:

\begin{itemize}
    \item \textbf{Síntesis de imágenes incondicional y condicional},
    incluyendo la generación texto-a-imagen a gran escala, viable
    gracias a los modelos de difusión latente que operan en un espacio
    comprimido en lugar de en el espacio de píxeles~\cite{rombach2022ldm}.
    \item \textbf{Problemas inversos en imagen}: super-resolución,
    \textit{inpainting}, \textit{deblurring} y \emph{denoising}, donde
    se busca recuperar una señal limpia a partir de una observación
    degradada.
    \item \textbf{Imagen médica}, que es el contexto de este trabajo:
    reconstrucción, reducción de ruido y síntesis de modalidades, con
    aplicaciones específicas a PET~\cite{scorepet2024,ddpmpet2023} que
    motivan directamente la contribución de esta memoria.
    \item \textbf{Otros dominios}: generación de audio y voz, de vídeo,
    de estructuras moleculares y de datos científicos, donde la misma
    receta de \emph{corromper y aprender a revertir} se traslada con
    pocas modificaciones.
\end{itemize}

La aplicación que vertebra este trabajo —la reconstrucción de PET de
baja dosis— se sitúa en la intersección entre los problemas inversos
en imagen y la imagen médica, y se desarrolla en el
\Cref{cap:pet}.

% ---------------------------------------------------------------------
\section{La tomografía por emisión de positrones}
\label{sec:pet-fisica}

Antes de revisar el estado del arte conviene fijar, brevemente, cómo
se forma una imagen PET y de dónde procede el ruido que este trabajo
intenta eliminar.

\subsection{Principio de funcionamiento}
\label{subsec:pet-funcionamiento}

Al paciente se le inyecta un \emph{radiotrazador}: una molécula
biológicamente activa —típicamente fluorodesoxiglucosa, un análogo de
la glucosa— marcada con un isótopo emisor de positrones de vida
corta, como el flúor-18. El trazador se distribuye por el torrente
sanguíneo y se acumula en los tejidos con mayor actividad metabólica:
el cerebro, el corazón, una lesión tumoral. Cada desintegración del
isótopo emite un \emph{positrón}, la antipartícula del electrón, que
recorre del orden de un milímetro de tejido antes de aniquilarse con
un electrón; la aniquilación produce dos fotones gamma de 511~keV
que salen en direcciones casi opuestas. Un anillo de detectores
alrededor del paciente registra los pares de fotones que llegan en
coincidencia temporal; cada coincidencia define una \emph{línea de
respuesta} sobre la que tuvo lugar la aniquilación. A partir de
millones de estas líneas, un algoritmo de reconstrucción tomográfica
—clásicamente retroproyección filtrada o métodos iterativos de
máxima verosimilitud tipo EM— estima la distribución espacial del
trazador. El resultado es un mapa funcional del metabolismo, no una
imagen anatómica: esa es la diferencia esencial del PET frente a los
rayos X o la TAC.

\subsection{Dosis, tiempo de adquisición y ruido}
\label{sec:pet-ruido}

La imagen PET es, en esencia, un recuento de eventos discretos, y por
ello su ruido es de naturaleza \textbf{Poisson}: la varianza del
recuento en cada elemento de medida es proporcional a su media, de
modo que la relación señal-ruido crece con la raíz del número de
eventos detectados. El número total de eventos es proporcional a la
dosis inyectada y al tiempo de adquisición. Reducir la dosis a un
veinteavo —el escenario de este trabajo— equivale a recortar el
recuento en ese mismo factor, con una degradación severa: en el
conjunto de datos empleado, la desviación típica del residuo entre la
versión de dosis baja y la de dosis completa, restringida al
tejido, es del orden del 31\,\% de la media de señal
(\Cref{sec:pet-datos}).

Un malentendido frecuente es que el modelo de ruido de Poisson del
PET sería incompatible con el ruido gaussiano artificial del proceso
de difusión. No es un problema: el proceso de difusión no pretende
imitar el ruido físico del instrumento, sino establecer una
correspondencia entre la distribución de imágenes de interés y una
distribución de referencia conocida, que se elige gaussiana por sus
propiedades matemáticas (\Cref{subsec:forward}). El ruido físico de
la observación entra en escena por otra vía: en el paradigma
supervisado, a través de los pares de entrenamiento que muestran al
modelo cómo es la degradación real; en el paradigma de guiado, a
través del término de verosimilitud de la \Cref{eq:dps-update}.

\subsection{La reconstrucción de dosis baja como problema inverso}
\label{subsec:pet-problema-inverso}

Formalmente, se dispone de una observación $\mathbf{y}$ (el escaneo
de dosis baja) y se busca la señal $\mathbf{x}$ (el escaneo de dosis
completa equivalente). La operación que liga ambas no es trivialmente
invertible: combina el re-muestreo Poisson del recuento, la reducción
del tiempo efectivo de adquisición y la etapa de reconstrucción
algorítmica del propio escáner, y destruye información. El problema
es por tanto \emph{mal planteado}, y toda solución se apoya —explícita
o implícitamente— en un prior sobre qué imágenes son plausibles. Los
dos paradigmas de la \Cref{subsec:difusion-condicional} son
exactamente las dos maneras de instanciar ese prior con un modelo de
difusión, y su comparación empírica sobre PET real es el objeto del
\Cref{cap:pet}.

% ---------------------------------------------------------------------
\section{Estado del arte en reconstrucción de PET con modelos de
difusión}
\label{sec:estado-del-arte}

Establecido el aparato conceptual y formal, esta sección recorre las
contribuciones más relevantes aplicadas al problema de la
reconstrucción de PET de dosis baja. La presentación es
\emph{cronológica}, de las propuestas más antiguas a las más
recientes, con la intención de mostrar cómo ha evolucionado el campo
y qué aporta cada trabajo respecto a los anteriores.

\paragraph{Octubre 2023 — DDPM como modelo de prior para
\textit{denoising} de PET~\cite{ddpmpet2023}.} La primera de las
contribuciones que se reseñan plantea, de manera todavía
exploratoria, la pregunta que vertebrará buena parte del campo
durante los años siguientes: ¿pueden los \emph{Denoising Diffusion
Probabilistic Models}, recién popularizados por el trabajo de Ho
et al.~\cite{ho2020ddpm}, utilizarse como modelo de aprendizaje
de la distribución de imágenes PET de buena calidad y, a partir
de ahí, reducir el ruido de escaneos degradados? El estudio
entrena un DDPM como modelo de distribución sobre escaneos PET y
compara varias estrategias derivadas para la tarea de
\textit{denoising}. El interés de este trabajo no es tanto el de
una propuesta cerrada cuanto el de una \emph{prueba de concepto}
que abre la puerta a las líneas posteriores: introduce los
ingredientes —un DDPM como prior aprendido, una forma de
condicionar o guiar el muestreo hacia consistencia con la
observación— sobre los que se construirán las propuestas
sofisticadas que siguen.

\paragraph{Enero 2024 — Score-based generative models aplicados a
PET~\cite{scorepet2024}.} Esta contribución traslada al dominio
del PET el marco de los \emph{score-based generative models},
formalizado por Song et al.~\cite{song2021sde} como una
unificación de los DDPM con la dinámica de Langevin sobre el
score. La reformulación tiene dos consecuencias prácticas
relevantes. La primera es \emph{teórica}: al expresar el proceso
de difusión como una ecuación diferencial estocástica continua,
se dispone de toda la maquinaria del cálculo estocástico —incluida
la ecuación inversa de Anderson— para razonar sobre cómo inyectar
la consistencia con la medición en el muestreo. La segunda es
\emph{operativa}: la flexibilidad para elegir muestreadores y para
incorporar términos de guiado permite atacar directamente el
problema inverso de la reconstrucción sin recurrir a un
entrenamiento condicional supervisado. Es, por tanto, el primer
trabajo que sitúa el guiado posterior como una alternativa
defendible al condicionamiento supervisado en PET.

\paragraph{Junio 2024 — DDPEM: integración del proceso EM con
DDPM~\cite{ddpem2024}.} Esta propuesta, denominada DDPEM, plantea
una hibridación interesante entre el proceso clásico de
reconstrucción \emph{Expectation Maximization} —el algoritmo de
referencia para la reconstrucción analítica de PET— y el muestreo
inverso de un DDPM. La motivación es físicamente coherente: la
operación de adquisición del PET admite una formulación
probabilística para la que EM proporciona un algoritmo de máxima
verosimilitud bien fundamentado, y combinarla con un prior
generativo aprendido por un DDPM permite incorporar conocimiento
sobre la estructura de las imágenes plausibles que el algoritmo
EM, por sí solo, no posee. El resultado es un esquema que
explícitamente reconoce el modelo físico del PET en lugar de
tratar la observación como una imagen genérica degradada. Esta
línea representa el extremo más \emph{model-based} del espectro
de soluciones basadas en difusión.

\paragraph{Junio 2024 — Revisión BJR\,AI: la difusión no supervisada
como estado del arte~\cite{bjrai2024}.} Apenas unas semanas
después, la revista \emph{British Journal of Radiology, Artificial
Intelligence} publica una revisión exhaustiva sobre el uso de
algoritmos de reconstrucción basados en modelos generativos
profundos para imagen médica que tiene un valor particular por dos
motivos. En primer lugar, \emph{establece} —tras una comparación
sistemática de la literatura disponible— que los algoritmos que
incorporan modelos de difusión \emph{no supervisados} constituyen
el estado del arte en tres tareas clave de imagen médica: la
resonancia magnética ultrarrápida, la TAC \textit{super-sparse-view}
y, de forma directamente relevante para este trabajo, el PET de
dosis baja. En segundo lugar, ofrece una introducción didáctica
muy cuidada tanto al planteamiento general de los problemas
inversos en imagen como a la maquinaria de los modelos de difusión,
acompañada de una discusión metodológica sobre los retos
específicos de cada modalidad (formulación del operador directo,
ruido característico, métricas adecuadas) y de un catálogo de los
campos de investigación abiertos. Más que una propuesta
metodológica, este trabajo funciona como \emph{punto de anclaje}
del estado del arte y como justificación para apostar por la línea
del prior no condicional con guiado de medición, que es uno de los
dos \textit{pipelines} desarrollados en esta memoria.

\paragraph{Junio 2025 — MAMBA para reconstrucción con
incertidumbre~\cite{mamba2025}.} Una vez asentada la difusión como
referencia para el dominio, la innovación se desplaza hacia la
\emph{arquitectura} de la red de \textit{denoising}. Este trabajo
sustituye la arquitectura tradicional de tipo U-Net o
\textit{transformer} por un modelo basado en el paradigma
\texttt{MAMBA}, una familia de modelos de espacio de estados que
combina eficiencia computacional —complejidad lineal en la
secuencia— con la capacidad de modelar dependencias de largo
alcance, particularmente atractiva para datos volumétricos en los
que la atención completa resulta prohibitiva. Además de proponer
una mejora arquitectónica, el trabajo incorpora una modelización
explícita de la \emph{incertidumbre} de la reconstrucción, un
ingrediente clínicamente relevante que tiende a estar ausente en
las propuestas anteriores. El estudio reclama el estado del arte
para PET de dosis baja, así como para resonancia magnética rápida
y TAC \textit{sparse-view}.

\paragraph{Diciembre 2025 — Task-Adaptive Transformer
(TAT)~\cite{tat2025}.} La contribución más reciente que se reseña
da otro giro arquitectónico: propone un \emph{Task-Adaptive
Transformer}, un único modelo que comparte representación entre
tareas de restauración heterogéneas y se especializa
dinámicamente en función de la tarea —síntesis de PET, reducción
de ruido en TAC, super-resolución en resonancia magnética—. La
filosofía es la de los \emph{modelos generalistas} de imagen
médica, un movimiento paralelo al de los grandes modelos
fundacionales en otros dominios. El estudio reclama el estado del
arte en las tres tareas que aborda. Para los propósitos de este
trabajo, su interés es relativo: aporta un techo competitivo de
referencia, pero queda fuera del alcance —y del presupuesto de
cómputo— de un TFM tratar de reproducir o superar un modelo
generalista de este tipo.

\paragraph{Lectura conjunta.} Vistas las contribuciones en orden
cronológico, el campo dibuja una trayectoria nítida: arranca en
2023 con DDPM básicos aplicados a PET~\cite{ddpmpet2023}, se
formaliza en 2024 mediante el marco de los \emph{score-based
generative models}~\cite{scorepet2024} y la hibridación con el
modelo físico del PET~\cite{ddpem2024}, se consolida ese mismo año
con una revisión que lo sitúa explícitamente como estado del
arte~\cite{bjrai2024}, y en 2025 se diversifica en dos direcciones
arquitectónicas —modelos de espacio de estados con cuantificación
de incertidumbre~\cite{mamba2025} y \textit{transformers}
generalistas multitarea~\cite{tat2025}— que comparten el sustrato
de difusión pero exploran maneras distintas de aprovecharlo. Esta
trayectoria justifica el diseño que se desarrolla en el
\Cref{cap:pet}: la elección de comparar un \textit{pipeline}
supervisado condicional —la formulación canónica más directa— con
un \textit{pipeline} de prior no condicional con guiado por
medición —la línea identificada como estado del arte por la
revisión BJR\,AI— es exactamente la decisión que el panorama
descrito sugiere.
